%=====================================================================
\section{Conclusion}
\label{sec:conclusion}

With a guideline and a handful of gold documents, the thing to learn is the procedure, not a model. Decision orchestration lets code own control flow and well-formedness and consults a model only for local decisions. Small models can then annotate a rich scheme, and the procedure can be calibrated, routed and corrected from a few documents through errors attributed to single nodes. The same explicitness shows the limit: with so few documents, the acceptance test, not the reflector, bounds how far learned edits generalize. \todo{second benchmark (E9) and human routing (E5)}

\section*{Limitations}
Our main evidence comes from one task in one language. The second benchmark (\S\ref{sec:res-second}) has not been run yet, and the English \umr runs, made with earlier versions of the procedure, are not reported. Test sets are small (99--157 sentences) and differ between stages of the work. Unifying them is planned (\exper{E1}). Most configurations were run once, and run-to-run noise (up to 0.6 \F for closed mode and 2.3 \F for a 9B model in fragment mode) is of the same order as some of the differences we report. The seeds were written by the authors, so their quality, and the effort it took to write them, affect every result. We have not measured that effort. Several of the models we use, including all three decision-model services and the frontier LLMs, are commercial APIs whose behavior may change; we release the cached answers to make runs reproducible \todo{check that the release includes them}. The routing results come from 20 sentences, and the ``human'' in routing will be simulated from gold, which overstates real annotators' accuracy. Finally, AnCast and our fine-grained judge depend on an automatic alignment, and the gold itself contains errors. One gold graph has unbalanced brackets and is skipped by the scorer.

\section*{Ethics Statement}
The data are publicly released Chinese \umr annotations. We collected no new data from people and no personal information beyond what the released documents contain. Our goal is to reduce the cost of producing the first annotations for a new scheme, which could reduce demand for annotation work. We see the method as a way to direct annotators' effort to the decisions that need human judgment. Running the experiments in this paper used commercial APIs, at a total cost we will report with the camera-ready version \todo{sum the API costs}.
